\documentclass[conference]{IEEEtran}
\IEEEoverridecommandlockouts
% Paquetes para español
\usepackage[spanish]{babel}
\usepackage[utf8]{inputenc}
% The preceding line is only needed to identify funding in the first footnote. If that is unneeded, please comment it out.
\usepackage{cite}
\usepackage{amsmath,amssymb,amsfonts}
\usepackage{algorithmic}
\usepackage{graphicx}
\usepackage{textcomp}
\usepackage{booktabs}
\usepackage{xcolor}
% Si falta una imagen, muestra un recuadro en lugar de detener la compilación
\let\oldincludegraphics\includegraphics
\renewcommand{\includegraphics}[2][]{\IfFileExists{#2}{\oldincludegraphics[#1]{#2}}{\fbox{\parbox{0.8\columnwidth}{\centering Imagen faltante:\\ \texttt{\detokenize{#2}}}}}}
\def\BibTeX{{\rm B\kern-.05em{\sc i\kern-.025em b}\kern-.08em
    T\kern-.1667em\lower.7ex\hbox{E}\kern-.125emX}}
\begin{document}

\title{Perceptrón}

\author{\IEEEauthorblockN{1\textsuperscript{st} Oscar Alberto Gallo García}
\IEEEauthorblockA{\textit{M. C. en Robótica e Inteligencia Artificial} \\
\textit{Centro Universitario de Ciencias Exactas e Ingenierías}\\
Guadalajara, México \\
oscar.gallo5770@alumnos.udg.mx}
}

\maketitle

\begin{abstract}
    Se trabajó el algoritmo simple del perceptrón desde cero, implementado para
    una neurona para clasificación lineal binaria y con dos neuronas para
    clasificación de cuatro clases.
    Se desarrollaron las funciones de entrenamiento y prediccion demostrando
    la capacidad de este algoritmo para resolver problemas linealmente separables.
\end{abstract}

\begin{IEEEkeywords}
Perceptrón, Neurona Artificial, Clasificacion lineal
\end{IEEEkeywords}

\section{Introducción}



\section{Desarrollo}

\subsection{Modelo del Sistema No Lineal}
El sistema de este proyecto se puede analizar a partir de las leyes de movimiento de Newton, dónde las fuerzas que actúan sobre el sistema se determinar a partir del esquemático del sistema (Figura \ref{fig:PICdiag}).\\
Es fácil analizar el sistema si lo separamos en dos subsistemas; el carro y el péndulo invertido.\\
Para el carro con desplazamiento lineal tenemos la suma de fuerzas \eqref{eq:Fcarro}
\begin{align}
    \label{eq:Fcarro}
    \sum F = 0 \rightarrow \quad F_c = H + M_c\ddot{x}_c+b_c\dot{x}_c
\end{align}
Dónde $F_c$ es la fuerza lineal que se ejerce sobre el carro, $H$ es la fuerza de reacción interna horizontal del pivote que conecta con el péndulo, $M_c$ es la masa de carro, $b_c$ es un coeficiente de fricción viscosa del carro y $x_c$ es el desplazamiento lineal del carro.
\begin{figure}[htbp]
\centering
\includegraphics[scale=0.1]{figures/PICdiag.jpg}
\caption{Esquemático del Sistema}
\label{fig:PICdiag}
\end{figure}
Ahora Podemos definir la suma de Torques en el péndulo.
\begin{align}
    \label{eq:Tpendulo}
    \sum \tau = 0\rightarrow \quad I\ddot{\alpha} = -Vl_p\sin(\alpha) -Hl_p\cos(\alpha) + b_p\dot{\alpha}
\end{align}
Dónde $V$ es la fuerza de reacción interna vertical del pivote que conecta con el péndulo, $I$ es el momento de inercia del péndulo respecto a su centro de gravedad, $l_p$ es la longitud del péndulo al centro de gravedad y $b_p$ es el coeficiente de fricción viscosa del péndulo.\\
Si aplicamos la segunda ley de Newton, podemos determinar el movimiento horizontal y vertical del péndulo \cite{b1}.
\begin{align}
    \label{eq:Fxpendulo}
    \sum F_x &= 0 \rightarrow \quad H = M_p\ddot{x}_p\\
    \label{eq:Fypendulo}
    \sum F_y &= 0 \rightarrow \quad V = M_pg + M_p\ddot{y}_p
\end{align}
Dónde $M_p$ es la masa del péndulo y $g$ es la constante de aceleración por efecto de la gravedad. Ahora determinamos la posición del péndulo en sus componentes $(x,y)$.
\begin{align}
    \label{eq:Pospendulo}
    x_p = x_c - l_p\sin(\alpha),\quad y_p = l_p\cos(\alpha)
\end{align}
Sustituyendo las coordenadas de \eqref{eq:Pospendulo} en \eqref{eq:Fxpendulo} y \eqref{eq:Fypendulo} tenemos:
\begin{align}
    H &= M_p\left( \ddot{x}_c - \frac{d^2}{dt^2}l_p\sin(\alpha) \right)\nonumber\\
    \label{eq:Fxpendulo2}
    H &= M_p\ddot{x}_c - M_pl_p\ddot{\alpha}\cos(\alpha) + M_pl_p\dot{\alpha}^2\sin(\alpha)\\
    V &= M_pg + M_p\left(\frac{d^2}{dt^2}l_p\cos(\alpha)\right)\nonumber\\
    \label{eq:Fypendulo2}
    V &= M_pg -M_pl_p\ddot{\alpha}\sin(\alpha) -M_pl_p\dot{\alpha}^2\cos(\alpha)
\end{align}
Ahora podemos sustituir \eqref{eq:Fxpendulo2} en \eqref{eq:Fcarro}
\begin{align}
    \label{eq:EcCarro}
    F_c = (M_p + M_c)\ddot{x}_c +M_pl_p\left(\dot{\alpha}^2\sin(\alpha)-\ddot{\alpha}\cos(\alpha) \right) + b_c\dot{x}_c
\end{align}
Y también podemos sustituir \eqref{eq:Fxpendulo2} y \eqref{eq:Fypendulo2} en \eqref{eq:Tpendulo}
\begin{align}
    \label{eq:EcPendulo}
    I\ddot{\alpha} =& M_pl_p\left( l_p\ddot{\alpha} - g\sin(\alpha) - \ddot{x}_c\cos(\alpha) \right) + b_p\dot{\alpha}
\end{align}
Las ecuaciones \eqref{eq:EcCarro} y \eqref{eq:EcPendulo} representan las ecuaciones de movimiento del sistema PIC (Péndulo Invertido sobre Carro).

\subsection{Espacio de Estados}
Tenemos que representar el sistema en espacio de estados de la forma \eqref{eq:SSForm}.
\begin{align}
    \label{eq:SSForm}
    \dot{\mathbf{x}}&=f(\mathbf{x})+g(\mathbf{x})u\\
    \mathbf{y}&=h(\mathbf{x})\nonumber
\end{align}
Comenzamos con definir el sistema en forma matricial
\begin{align*}
    \begin{bmatrix}
        M_pl^2_p - I & -M_pl_p\cos(\alpha)\\
        -M_pl_p\cos(\alpha) & M_p + M_c
    \end{bmatrix}
    \begin{bmatrix}
        \ddot{\alpha} \\ \ddot{x}
    \end{bmatrix}
    =\\
    \begin{bmatrix}
        M_pl_pg\sin(\alpha) - b_p\dot{\alpha}\\
        F_c -M_pl_p\dot{\alpha}^2\sin(\alpha)-b_c\dot{x}_c
    \end{bmatrix}
\end{align*}
Despejamos el vector de aceleraciones
\begin{align*}
    \begin{bmatrix}
        \ddot{\alpha} \\ \ddot{x}
    \end{bmatrix}
    =
    \begin{bmatrix}
        M_pl^2_p - I & -M_pl_p\cos(\alpha)\\
        -M_pl_p\cos(\alpha) & M_p + M_c
    \end{bmatrix}^{-1}\\
    \begin{bmatrix}
        M_pl_pg\sin(\alpha) - b_p\dot{\alpha}\\
        F_c -M_pl_p\dot{\alpha}^2\sin(\alpha)-b_c\dot{x}_c
    \end{bmatrix}
\end{align*} 
Al calcular la inversa y multiplicar las matrices tenemos dos ecuaciones para cada aceleración. Aprovechamos este paso para hacer el cambio de variables a variables de estado:
\begin{align*}
    x_1=x_c,\quad x_2=\dot{x}_c,\quad x_3=\alpha,\quad x_4=\dot{\alpha}
\end{align*}
Y el espacio de estados queda expresado mediante \eqref{eq:StateSpace}
\begin{align}
    \label{eq:StateSpace}
    \begin{bmatrix}
        \dot{x}_1\\ \dot{x}_2\\ \dot{x}_3\\ \dot{x}_4
    \end{bmatrix}=
    \begin{bmatrix}
        x_2\\ f_2(\mathbf{x})\\ x_4 \\ f_4(\mathbf{x)}
    \end{bmatrix}+
    \begin{bmatrix}
        0\\g_2(\mathbf{x})\\ 0\\ g_4(\mathbf{x})
    \end{bmatrix}F_c
\end{align}
Dónde
\begin{align}
     f_2(\mathbf{x}) =&
     [-(I+M_p l_p^2)\,b_c\,x_2 - (M_p^2 l_p^3 + M_p l_p I)\,x_4^2 \sin(x_3)\nonumber\\
     \label{eq:f2}
      &- M_p l_p \cos(x_3)\,b_p\,x_4
        + M_p^2 l_p^2 g \cos(x_3)\sin(x_3)]/D(\mathbf{x})\\
    f_4(\mathbf{x}) =&
     [(M_p+M_c)M_p l_p g \sin(x_3) - (M_p+M_c)\,b_p\,x_4\nonumber\\
     \label{eq:f4}
      &- M_p^2 l_p^2 x_4^2 \sin(x_3)\cos(x_3)
    - M_p l_p b_c \cos(x_3)\,x_2]/D(\mathbf{x})\\
    \label{eq:g2}
    g_2(\mathbf{x})=& \frac{I + M_p l_p^2}{D(\mathbf{x})}\\
    \label{eq:g4}
    g_4(\mathbf{x})=& \frac{M_p l_p \cos(x_3)}{D(\mathbf{x})}\\
    D(\mathbf{x})=& (M_p+M_c)I+M_pM_cl_p^2 +M_p^2l_p^2\sin^2(x_3)\nonumber
\end{align}
\subsection{Puntos de Equilibrio}
Para obtener los puntos de equilibrio tenemos que evaluar $\dot{\mathbf{x}}=\mathbf{0}$
\begin{align*}
    \dot{x}_1=0\rightarrow\quad x_2=0\\
    \dot{x}_3=0\rightarrow\quad x_4=0\\
\end{align*}
\begin{align}
    &\dot{x}_2=0\rightarrow\quad f_2(\mathbf{x})+g_2(\mathbf{x})F_c=0\nonumber\\
    &(I+M_pl_p^2)F_c+M_p^2l_p^2g\cos(x_3)\sin(x_3)=0\nonumber\\
    \label{eq:F1Equilibrium}
    &F_c=-\frac{M_p^2l_p^2g\cos(x_3)\sin(x_3)}{I+M_pl_p^2}
\end{align}
\begin{align}
    &\dot{x}_4=0\rightarrow\quad f_4(\mathbf{x})+g_4(\mathbf{x})F_c=0\nonumber\\
    \label{eq:F2Equilibrium}
    &(M_p+M_c)g\sin(x_3)+F_c\cos(x_3)=0
\end{align}
Sustituyendo \eqref{eq:F1Equilibrium} en \eqref{eq:F2Equilibrium}
\begin{align*}
    g\sin(x_3)\left(M_p+M_c-\frac{M_p^2l_p^2\cos^2(x_3)}{I+M_pl_p^2}\right)=0
\end{align*}
De aquí podemos tener dos factores que podrían ser 0. En el primer factor
\begin{align}
    \label{eq:FirstFactor}
    g\sin(x_3)=0\rightarrow\quad x_3=k\pi\ |\ k\in\mathbb{Z}
\end{align}
En el segundo factor
\begin{align}
    \label{eq:SecondFactor}
    \cos^2(x_3)=\frac{(M_p+M_c)(I+M_pl_p^2)}{M_p^2l_p^2}
\end{align}
Sabemos que
\begin{align*}
    0\leq \cos^2(x_3)\leq 1,\forall x_3\in\mathbb{R}
\end{align*}
Por lo tanto, para que este segundo factor sea 0, debe cumplirse:
\begin{align}
    \label{eq:FirstIneq}
    (M_p+M_c)(I+M_pl_p^2)\leq M_p^2l_p^2
\end{align}
Pero observamos que
\begin{align*}
    M_p+M_c>0,\quad I+M_pl_p^2>M_pl_p^2
\end{align*}
Entonces
\begin{align}
    \label{eq:SecondIneq}
    (M_p+M_c)(I+M_pl_p^2)> M_p(M_pl_p^2)
\end{align}
Las desigualdades \eqref{eq:FirstIneq} y \eqref{eq:SecondIneq} se contradicen, esto implica que \eqref{eq:SecondFactor} no se puede cumplir para parámetros físicos reales. Solo el factor \eqref{eq:FirstFactor} genera el 0 en la multiplicación.\\
Por último sustituimos $x_3=k\pi$ en \eqref{eq:F1Equilibrium}
\begin{align*}
    F_c=0
\end{align*}
Los puntos de equilibrio encontrados son:
\begin{align*}
    \mathbf{x}^*=
    \begin{bmatrix}
        x_1^* & 0 & k\pi & 0
    \end{bmatrix}^\top\ | \ k\in\mathbb{Z}
\end{align*}
Como se puede observar, existen infinitos puntos de equilibrio, pero se pueden simplificar en $x_3^*$ a los equilibrios naturales del péndulo (colgando e invertido)
\begin{align}
    \label{eq:Equilibrium}
    \mathbf{x}_1^*=
    \begin{bmatrix}
         x_1^* & 0 & 0 & 0
    \end{bmatrix}^\top,\quad
    \mathbf{x}_2^*=
    \begin{bmatrix}
         x_1^* & 0 & \pi & 0
    \end{bmatrix}^\top,\quad F_c=0
\end{align}
\subsection{Análisis de Puntos de Equilibrio}
Para determinar el comportamiento y estabilidad de los puntos de equilibrio, se escoge linealizar el sistema al rededor de dichos puntos. Para ello se comienza con el cálculo del Jacobiano:
\begin{align}
    \label{eq:Jacobian}
    \left. \begin{bmatrix}
        0 & 1 & 0 & 0\\
        0 & \frac{\partial f_2}{\partial x_2} & \frac{\partial f_2}{\partial x_3} & \frac{\partial f_2}{\partial x_4}\\
        0 & 0 & 0 & 1\\
        0 & \frac{\partial f_4}{\partial x_2} & \frac{\partial f_4}{\partial x_3} & \frac{\partial f_4}{\partial x_4}\\
    \end{bmatrix}\right|_{\mathbf{x}=\mathbf{x}^*} =
    \begin{bmatrix}
        0 & 1 & 0 & 0\\
        0 & a_{22} & a_{23} & c_3a_{24}\\
        0 & 0 & 0 & 1\\
        0 & c_3a_{42} & c_3a_{43} & a_{44}\\
    \end{bmatrix}
\end{align}
Con
\begin{align*}
    a_{22} &= -\frac{(I+M_pl_p^2)b_c}{D_{\mathbf{x}^*}},
    &&a_{24} = -\frac{M_pl_pb_p}{D_{\mathbf{x}^*}}\\
    a_{23} &= \frac{M_p^2l_p^2g}{D_{\mathbf{x}^*}}\\
    a_{42} &= -\frac{M_pl_pb_c}{D_{\mathbf{x}^*}},
    &&a_{44} = -\frac{(M_p+M_c)b_p}{D_{\mathbf{x}^*}}
\end{align*}
\begin{align*}
    a_{43} &= \frac{(M_p+M_c)M_pl_pg}{D_{\mathbf{x}^*}},\quad c_3=\cos(x_3)
\end{align*}
Dónde
\begin{align*}
    D_{\mathbf{x}^*}=(M_p+M_c)I+M_pM_cl_p^2
\end{align*}
Evaluamos la matriz en $\mathbf{x}_1^*$
\begin{align}
    \label{eq:Aeq1}
    A_1=
    \begin{bmatrix}
        0 & 1 & 0 & 0\\ 0 & a_{22} & a_{23} & a_{24}\\ 0 & 0 & 0 & 1\\
        0 & a_{42} & a_{43} & a_{44}
    \end{bmatrix}
\end{align}
Y en $\mathbf{x}_2^*$
\begin{align}
    \label{eq:Aeq2}
    A_2=
    \begin{bmatrix}
        0 & 1 & 0 & 0\\ 0 & a_{22} & a_{23} & -a_{24}\\ 0 & 0 & 0 & 1\\
        0 & -a_{42} & -a_{43} & a_{44}
    \end{bmatrix}
\end{align}
Para analizar el punto de equilibrio, tenemos que calcular los valores propios de la matriz $A$ correspondiente. Para simplificar los cálculos, se sustituyen los valores reales de los parámetros físicos del experimento, descritos en el Cuadro \ref{tab:params}.\\
Para $A_1$
\begin{align*}
    \lambda_1=0,\quad \lambda_2=-8.2350,\quad \lambda_3=-4.6498,\quad \lambda_4=6.5
\end{align*}
Se puede ver que existe al menos un valor propio real positivo, por lo tanto $\mathbf{x}_1^*$ es inestable.\\
Para $A_2$
\begin{align*}
    \lambda_1=0,\quad \lambda_2=-5.4542,\quad \lambda_3=-0.4695\pm6.7460i
\end{align*}
Existen valores propios negativos con parte imaginaria conjugada, por lo tanto $\mathbf{x}_2^*$ es marginalmente estable.

\subsection{Control}
Para el control no lineal se escoge realizar una linealización entrada-salida, que que busca simplificar la relación entre la entrada de control ($F_c=u$) y una salida específica del sistema. En nuestro caso definimos como salida la posición del péndulo, es decir, $x_3$.
\begin{align*}
    y=h(\mathbf{x})=x_3
\end{align*}
Comenzamos calculando la derivada de $y$ repetidamente hasta que la entrada $u$ aparezca explícitamente con un coeficiente no nulo. Esté numero de veces que derivamos define el grado relativo $r$
\begin{align*}
    \dot{y}&=\frac{\partial h(\mathbf{x})}{\partial \mathbf{x}}[f(\mathbf{x})+g(\mathbf{x})u]=x_4=\psi_2(\mathbf{x})\\
    \ddot{y}&=\frac{\partial \psi_2(\mathbf{x})}{\partial \mathbf{x}}[f(\mathbf{x})+g(\mathbf{x})u]=f_4(\mathbf{x})+g_4(\mathbf{x})u
\end{align*}
en la derivada 2 se encuentra $u$, por lo tanto $r=2<n$.\\
Definimos el sistema entrada-salida con $z_1=y,\ z_2=\dot{y}$
\begin{align}
    \label{eq:InOutSis}
    \dot{z_1} &= z_2\\
    \dot{z_2} &= f_4(\mathbf{x})+g_4(\mathbf{x})u\nonumber
\end{align}
Definimos $u$ para linealizar el sistema entrada-salida
\begin{align}
    \label{eq:UnonLinear}
    u=\frac{v-f_4(\mathbf{x})}{g_4(\mathbf{x})}
\end{align}
Al sustituir \eqref{eq:UnonLinear} en \eqref{eq:InOutSis}, obtenemos el sistema lineal entrada-salida en lazo cerrado
\begin{align}
    \label{eq:InOutLinSis}
    \dot{z}_1 &= z_2\\
    \dot{z}_2 &= v\nonumber
\end{align}
Nótese que en $g_4(\mathbf{x})$ tenemos un $\cos(x_3)$, por lo tanto, la linealización con $u$ solamente es válida si:
\begin{align*}
    -\frac{\pi}{2}\leq x_3\leq \frac{\pi}{2}
\end{align*}
Una vez linealizado el sistema entrada-salida, se puede diseñar $v$ con asignación de polos.
\begin{align*}
    v=-k_1y-k_2\dot{y}
\end{align*}
Entonces
\begin{align*}
    \dot{\mathbf{z}}= \mathbf{A}_z\mathbf{z}=
    \begin{bmatrix}
        0 & 1\\ -k_1 & -k_2
    \end{bmatrix}\mathbf{z}
\end{align*}
El polinomio característico es
\begin{align*}
    p_z(\lambda)=\det(\lambda\mathbf{I}-\mathbf{A}_z)=\lambda^2+k_2\lambda+k_1=0
\end{align*}
Si lo parametrizamos como un sistema de segundo orden
\begin{align*}
    p_z(\lambda)=\lambda^2+2\zeta\omega_n\lambda+\omega_n^2=0,\quad k_1=\omega_n^2,\ k_2=2\zeta\omega_n
\end{align*}
Escogemos una frecuencia $\omega_n<\sqrt{(M_pl_pg)/I_p}$. Sustituyendo los valores parámetros físicos tenemos  $\omega_n=6$ y $\zeta=0.8$. Esto nos da las ganancias $k_1=36$ y $k_2=9.6$.
\subsection{Simulación}


El sistema consiste en un péndulo rígido de masa $M_p$ articulado sobre un carro $M_c$ que se desplaza horizontalmente. La Tabla~\ref{tab:params} resume los parámetros del hardware empleado.

\begin{table}[htbp]
\caption{Parámetros del sistema Quanser IP01 + SIP}
\label{tab:params}
\begin{center}
\begin{tabular}{lcc}
\toprule
\textbf{Parámetro} & \textbf{Símbolo} & \textbf{Valor} \\
\midrule
Masa del péndulo        & $M_p$   & $0.127$~kg \\
Longitud total          & $L_p$   & $0.3365$~m \\
Dist.\ al c.d.m.        & $l_p$   & $0.1778$~m \\
Momento de inercia      & $I_p$   & $1.20\!\times\!10^{-3}$~kg\,m$^2$ \\
Amort.\ péndulo         & $B_p$   & $0.0024$~N\,m\,s/rad \\
Masa del carro          & $M_c$   & $0.7031$~kg \\
Amort.\ equivalente     & $B_{eq}$& $4.3$~N\,s/m \\
Gravedad                & $g$     & $9.81$~m/s$^2$ \\
\bottomrule
\end{tabular}
\end{center}
\end{table}

\begin{figure}[htbp]
\centering
\includegraphics[scale=0.15]{figures/image_1.png}
\caption{Simulación: Posición del carro}
\label{fig:PICSimpos}
\end{figure}

\begin{figure}[htbp]
\centering
\includegraphics[scale=0.15]{figures/image_2.png}
\caption{Simulación: ángulo del péndulo}
\label{fig:PICSimdeg}
\end{figure}




\section{Resultados}

\begin{figure}[htbp]
\centering
\includegraphics[scale=0.5]{figures/3.png}
\caption{Sistema experimental}
\label{fig:PICexp}
\end{figure}

\begin{figure}[htbp]
\centering
\includegraphics[scale=0.2]{figures/image_3.png}
\caption{Sistema experimental: Posición del carro}
\label{fig:PICpos}
\end{figure}

\begin{figure}[htbp]
\centering
\includegraphics[scale=0.13]{figures/Angulo.jpg}
\caption{Sistema experimental: ángulo del péndulo}
\label{fig:PICdeg}
\end{figure}

\section{Conclusiones}
Aunque la condición teórica de la linealización requiere únicamente que $\cos(x_3)\neq0$, lo cual se cumple para $|x_3|<\pi/2$, en la práctica la validez del modelo se restringe a un entorno pequeño alrededor del punto de equilibrio.
Esto se debe a que, conforme $x_3$ se aproxima a $\pm\pi/2$, el término $g_4(\mathbf{x})$ tiende a cero, provocando que la ley de control requiera valores excesivamente grandes de la entrada $u$, lo cual no es físicamente realizable.
Además, para ángulos grandes, los efectos no lineales se vuelven dominantes, deteriorando la precisión de la linealización. Por esta razón, se considera un rango práctico de operación de aproximadamente $|x_3|<10\sim15\ \deg$.

\begin{thebibliography}{00}

\bibitem{b1} H. K. Khalil, Nonlinear Systems, 3rd ed. Upper Saddle River, NJ: Prentice Hall, 2002.

\bibitem{b2} J.-J. E. Slotine and W. Li, Applied Nonlinear Control. Englewood Cliffs, NJ: Prentice Hall, 1991.

\bibitem{b3} A. Isidori, Nonlinear Control Systems, 3rd ed. London, U.K.: Springer-Verlag, 1995.

\bibitem{b4} K. Ogata, Modern Control Engineering, 5th ed. Upper Saddle River, NJ: Prentice Hall, 2010.

\bibitem{b5} F. L. Lewis, D. Vrabie, and V. L. Syrmos, Optimal Control, 3rd ed. Hoboken, NJ: Wiley, 2012.

\end{thebibliography}

\end{document}
